What can division tell us about the primes? Why is a number a square modulo one prime but not another? How can a good rational approximation solve an equation in integers?

Begin with Euclid's algorithm. Follow it to B\'ezout's identity, unique prime factorisation and congruences. Count divisors with M\"obius inversion, then study primitive roots and quadratic reciprocity. Decide which numbers are sums of two squares; solve integer equations by descent and Pell's equation by continued fractions; and explore Gaussian and Eisenstein integers, characters and sieves.

Concrete questions lead to proofs, worked calculations, mathematical diagrams and all 144 graded exercises, each with a hint and a complete solution. The methods have limits: unique factorisation can fail in a larger ring, and an Euler product is not yet a proof of the prime number theorem. The groups, rings and characters encountered here lead directly to Volume~3, \emph{Symmetry}; ideal theory and deeper prime distribution come later.
